% Originally: content/week-10.tex, \section{Flux and the divergence theorem} (Corsin Week 10)

\section{Flux and the divergence theorem}
\label{sec:flux_divergence_theorem}

% Source: Corsin Nick/Class Notes/Week 10.pdf, p. 7
\begin{definition}[Flux]
\label{def:flux}
Suppose $F : U \to \mathbb{R}^n$ is a $C^1$ vector field, $U \subseteq \mathbb{R}^n$ open, and
$\Omega \subseteq U$ is a $C^1$ bounded domain. The \newterm{flux} \germanfor{Fluss} of $F$
through $\partial\Omega$ is defined as
\[\int_{\partial\Omega} \langle F,\nu\rangle\,d\vol_{n-1} \qquad \text{(analysis notation)}
  \quad = \quad \int_{\partial\Omega} \vec F(x)\cdot d\vec A \qquad \text{(physics notation)}\]
\[= \int_B \big\langle F\circ\phi(x),\ \nu\circ\phi(x)\big\rangle
   \sqrt{\det D\phi(x)\transp D\phi(x)}\,dx \qquad \text{(explicit)},\]
where $\nu : \partial\Omega \to \mathbb{S}^{n-1}$ is the exterior unit normal, i.e.\ $\nu(y)$
is the outward-facing unit normal of $\partial\Omega$ at $y \in \partial\Omega$, and
$\phi : B \to \partial\Omega$ is a (local) parametrization of $\partial\Omega$. In general,
multiple integrals in the explicit formulation may be necessary.
\end{definition}

% Source: Corsin Nick/Class Notes/Week 10.pdf, p. 8
\begin{proposition}[Flux in $\mathbb{R}^3$]
\label{prop:flux_r3}
Let $\Omega \subseteq \mathbb{R}^3$ be a bounded $C^1$ domain and $F \in
C^1(\overline\Omega,\mathbb{R}^3)$ a vector field, and let $\phi : B \to \partial\Omega
\setminus N$ be a parametrization of $\partial\Omega\setminus N$, where $N \subseteq
\partial\Omega$ is a null set, $B \subseteq \mathbb{R}^2$. Then:
\begin{enumerate}[label=\textbf{\arabic*.}]
  \item $\displaystyle\nu(x) := \frac{\partial_1\phi(x)\times\partial_2\phi(x)}
    {\lvert\partial_1\phi(x)\times\partial_2\phi(x)\rvert}$ is an (interior or exterior) Gauss
    map (more precisely, $\nu\circ\phi^{-1}$ is a Gauss map, i.e.\
    $\nu(x) \in N_{\phi(x)}(\partial\Omega\setminus N)$).
  \item The Gram determinant is given by
    $\sqrt{\det D\phi(x)\transp D\phi(x)} = \lvert\partial_1\phi(x)\times\partial_2\phi(x)
    \rvert$.
  \item Therefore,
    \[\int_{\partial\Omega} \vec F\cdot d\vec A = \pm\int_B \big\langle F\circ\phi(x),\
      \partial_1\phi(x)\times\partial_2\phi(x)\big\rangle\,dx,\]
    where the sign depends on whether $\nu(x)$ is interior ($-$) or exterior ($+$).
\end{enumerate}
\end{proposition}

% Source: Corsin Nick/Class Notes/Week 10.pdf, pp. 9--10
\begin{example}[Flux through the unit sphere]
\label{ex:flux_unit_sphere}
Compute the flux of $F(x,y,z) := (xz,\,z,\,y)$ through the surface of $\mathbb{S}^2$ with
interior $B_1(0)$.

We use spherical coordinates with $r\equiv1$:
\[\phi : (0,2\pi)\times(0,\pi) \to \mathbb{R}^3, \qquad (\varphi,\theta) \mapsto
  \begin{pmatrix}\sin\theta\cos\varphi \\ \sin\theta\sin\varphi \\ \cos\theta\end{pmatrix}.\]

\textbf{Normal vector.}
\[\partial_\varphi\phi = \begin{pmatrix}-\sin\theta\sin\varphi\\ \sin\theta\cos\varphi\\
  0\end{pmatrix}, \qquad \partial_\theta\phi = \begin{pmatrix}\cos\theta\cos\varphi\\
  \cos\theta\sin\varphi\\ -\sin\theta\end{pmatrix}
  \quad \Longrightarrow \quad \partial_\varphi\phi\times\partial_\theta\phi
  = \begin{pmatrix}-\sin^2\theta\cos\varphi\\ -\sin^2\theta\sin\varphi\\
  -\sin\theta\cos\theta\end{pmatrix}.\]

\textbf{Sanity check: does this make sense, and is it exterior or interior?} We have
\[\partial_\varphi\phi\times\partial_\theta\phi = -\sin\theta\ \phi(\varphi,\theta)
  = -\sin\theta\,(x,y,z),\]
and the exterior normal of $\mathbb{S}^2$ is $+(x,y,z)$, so $\partial_\varphi\phi\times
\partial_\theta\phi$ is an \emph{interior} normal (for $\theta\in(0,\pi)$, $\sin\theta>0$).
So we multiply by $-1$ to switch the order in the cross product, matching the exterior sign
in \cref{prop:flux_r3}.

\textbf{The integral.}
\[
\begin{aligned}
\int_{\mathbb{S}^2} \vec F\cdot d\vec A
  &= \int_0^{2\pi}\!\!\int_0^\pi \begin{pmatrix}xz\\ z\\ y\end{pmatrix}\cdot
     \begin{pmatrix}\sin\theta\cos\varphi\\ \sin\theta\sin\varphi\\ \cos\theta\end{pmatrix}
     \sin\theta\,d\theta\,d\varphi \\
  &= \int_0^{2\pi}\!\!\int_0^\pi \Big[\sin^3\theta\cos\theta\cos^2\varphi
     + \sin^2\theta\cos\theta\sin\varphi + \sin^2\theta\cos\theta\sin\varphi\Big]\,d\theta\,
     d\varphi \\
  &= \pi\int_0^\pi \sin^3\theta\cos\theta\,d\theta
   = \pi\left[\frac{\sin^4\theta}{4}\right]_0^\pi = 0,
\end{aligned}
\]
using $\int_0^{2\pi}\cos^2\varphi\,d\varphi = \pi$ and $\int_0^{2\pi}\sin\varphi\,d\varphi=0$
for the middle two terms.
\end{example}

% Source: Corsin Nick/Class Notes/Week 10.pdf, p. 11
\begin{theorem}[Gauss]
\label{thm:gauss_divergence}
Let $\Omega \subseteq \mathbb{R}^n$ be a bounded $C^1$ domain and $F : \overline\Omega \to
\mathbb{R}^n$ a $C^1$ vector field. Then
\[\int_\Omega \divg F\,d\vol_n = \int_{\partial\Omega} \langle F,\nu\rangle\,d\vol_{n-1},
  \qquad \text{or, in physics notation,} \qquad
  \int_\Omega \vec\nabla\cdot\vec F\,dx = \int_{\partial\Omega} \vec F\cdot d\vec A.\]
\end{theorem}

\begin{ainote}
Corsin dates the theorem to Lagrange (discovered 1762) and Ostrogradsky (proven 1828), and
adds, next to the physics form, an unusually enthusiastic \qt{how cool is this!} The remark is
earned: a volume integral of a derivative is reduced to a boundary integral of the original
field, which is the shape of the fundamental theorem of calculus.

% Supplement: Sascha Brack/Class Notes/Week_10_Notes_Monday.pdf, p. 13
% Revised 2026-08-09, two faults found by rendering it:
%   * The label was placed at (-1.5,-1), which is ON the boundary, not inside it: the ellipse
%     is x^2/4 + y^2/2.25 = 1, so at x = -1.5 it passes through y = -1.5*sqrt(1-2.25/4)
%     = -0.992. The ultra-thick curve struck straight through the glyph. Moved to (-1.5,0.45),
%     where the boundary is at |y| = 0.992, so the label sits well inside and clear of the
%     nearest source at (-0.5,0.5), whose spokes reach only to x = -0.9.
%   * The sinks were drawn `<-` from the centre outward, putting all five arrowheads on the
%     centre dot, exactly the blob that spoiled the sink/source figure in \cref{ch:vector_calculus}.
%     Sources and sinks now both run between radius 0.12 and 0.40, outward and inward
%     respectively, so the two have the same footprint and differ only in direction.
\begin{center}
\begin{tikzpicture}[>=Latex, scale=1.5]
  % Domain Omega
  \draw[ultra thick] (0,0) ellipse (2cm and 1.5cm);
  \node at (-1.5,0.45) {$\Omega$};
  \node[right] at (2,-0.5) {$\partial\Omega$};

  % Internal sources (red): rays run OUTWARD, from radius 0.12 to 0.40.
  \foreach \x/\y in {-0.5/0.5, 0.8/-0.2, 0/-0.6, -1/-0.3} {
    \fill[red] (\x,\y) circle (1.5pt);
    \foreach \angle in {0,72,...,288} {
      \draw[red, thick, ->] (\x,\y) ++(\angle:0.12) -- ++(\angle:0.28);
    }
  }

  % Internal sinks (blue): the same rays run INWARD, from radius 0.40 to 0.12.
  \foreach \x/\y in {1/0.6, 0.2/0.8, -0.2/-1.1, 1.2/-0.8} {
    \fill[blue] (\x,\y) circle (1.5pt);
    \foreach \angle in {0,72,...,288} {
      \draw[blue, thick, ->] (\x,\y) ++(\angle:0.40) -- ++({\angle+180}:0.28);
    }
  }

  % Boundary flux: outward where the field leaves, inward where it enters. The `<-` form puts
  % the head on the boundary point itself, so these do not blob.
  \foreach \angle in {0,45,90,135,180,225,270} {
    \draw[red, thick, ->] ({2*cos(\angle)},{1.5*sin(\angle)}) -- ++(\angle:0.4);
  }
  \foreach \angle in {30,150,315} {
    \draw[blue, thick, <-] ({2*cos(\angle)},{1.5*sin(\angle)}) -- ++(\angle:0.4);
  }
\end{tikzpicture}
\end{center}
\end{ainote}

\begin{example}[Recomputing \cref{ex:flux_unit_sphere} via the divergence theorem]
\label{ex:flux_unit_sphere_via_divergence}
With $\Omega = B_1(0)$ and $F(x,y,z) = (xz,z,y)$ as before, $\divg F = \partial_x(xz) +
\partial_y(z) + \partial_z(y) = z + 0 + 0 = z$. So, in spherical coordinates ($z =
r\cos\theta$),
\[
\int_{\partial\Omega} \vec F\cdot d\vec A = \int_\Omega \divg F\,dx = \int_\Omega z\,dx
  = \int_0^1\!\!\int_0^{2\pi}\!\!\int_0^\pi \underbrace{\cos\theta\sin\theta}_{=\,\frac12\sin(2\theta),\
    \pi\text{-periodic}} r^3\,d\theta\,d\varphi\,dr = 0,
\]
since $\int_0^\pi \sin(2\theta)\,d\theta = 0$. This confirms \cref{ex:flux_unit_sphere} with none
of its cross-product bookkeeping.
\end{example}

\begin{ainote}
Corsin writes the divergence as $\divg F = z$ directly without the intermediate
$\partial_x(xz)+\partial_z(z)+\partial_y(y)$ step; spelled out above for clarity. Worth
noting how much the two computations of the same flux differ in character: the direct
definition needs an explicit Gauss map, an orientation sanity check, and three cross terms
in the integrand, while the divergence theorem needs only $\divg F$ and a coordinate change
already used in \cref{ex:paraboloid_patch_area}. This is the theorem's main
practical use, and it is not conceptual elegance but the plain fact that $\divg F$ is often far
simpler than $F$ restricted to a curved boundary.
\end{ainote}

% Supplement: Sascha Brack/Class Notes/Week_10_Notes_Monday.pdf, p. 14
\begin{example}[Computing flux using the divergence theorem]
\label{ex:flux_divergence_cube}
Let $\Omega := [-1,1]^3 \subset \mathbb{R}^3$ and $F(x,y,z) := (x^2, y^2, z^2)\transp$.
We compute the flux of $F$ out of $\Omega$ using the divergence theorem.
First, we find the divergence of $F$:
\[\divg F(x,y,z) = \partial_x(x^2) + \partial_y(y^2) + \partial_z(z^2) = 2x + 2y + 2z.\]
Then, integrating over $\Omega$:
\begin{align*}
\int_{\partial\Omega} \langle F(p),\ \nu(p)\rangle \mathop{d\text{vol}}_2(p) &= \int_\Omega \divg F(x,y,z) \mathop{d\text{vol}}_3(x,y,z) \\
&= \int_{-1}^1 \int_{-1}^1 \int_{-1}^1 (2x + 2y + 2z) \,dx\,dy\,dz.
\end{align*}
By symmetry, the integral of $2x$ over $[-1,1]$ is $0$, and similarly for $y$ and $z$. Therefore, the total flux is $0$.
\end{example}

% Source: Jérôme Paschoud/Notizen/Woche 10.2 Divergenzsatz.pdf
\begin{example}[Flux out of a cylinder]
\label{ex:divergence_cylinder}
Let $\Omega := \{(x,y,z) \in \mathbb{R}^3 \mid x^2 + y^2 \leq R^2,\ 0 \leq z \leq H\}$ be a cylinder of radius $R$ and height $H$. Consider the vector field $F(x,y,z) := (x^3 + \sin(z), y^3 + xz, z^3 + xy)\transp$.
We compute the outward flux of $F$ through $\partial\Omega$ using the divergence theorem.
First, we compute the divergence of $F$:
\[\divg F(x,y,z) = 3x^2 + 3y^2 + 3z^2 = 3(x^2 + y^2 + z^2).\]
By the divergence theorem, the flux is the integral of the divergence over $\Omega$. We compute this using cylindrical coordinates $(r, \theta, z)$:
\begin{align*}
\int_{\partial\Omega} \langle F(p),\ \nu(p)\rangle \mathop{d\text{vol}}_2(p) &= \int_\Omega 3(x^2 + y^2 + z^2) \mathop{d\text{vol}}_3(x,y,z) \\
&= \int_0^H \int_0^{2\pi} \int_0^R 3(r^2 + z^2) r \,dr\,d\theta\,dz \\
&= 2\pi \int_0^H \left[ \frac{3}{4}r^4 + \frac{3}{2}r^2 z^2 \right]_0^R \,dz \\
&= 2\pi \int_0^H \left( \frac{3}{4}R^4 + \frac{3}{2}R^2 z^2 \right) \,dz \\
&= 2\pi \left[ \frac{3}{4}R^4 z + \frac{1}{2}R^2 z^3 \right]_0^H \\
&= 2\pi \left( \frac{3}{4}R^4 H + \frac{1}{2}R^2 H^3 \right) = \frac{\pi}{2}R^2 H (3R^2 + 2H^2).
\end{align*}
\end{example}

% Supplement: Analysis_II_Script_v1.pdf, p. 145
\begin{example}[Archimedes' principle]
\label{ex:archimedes_principle}
A solid $\Omega \subset \mathbb{R}^3$ (a bounded $C^1$ domain) is fully submerged in water, with
the coordinates chosen so that $x_3=0$ is the water's surface. At depth $-x_3$ the water
pressure is $P(x) := -\rho gx_3$, where $\rho$ is the density of water and $g$ the gravitational
acceleration. The force the water exerts on a surface element near $q \in \partial\Omega$ is
proportional to the pressure and points inward, along $-\nu(q)$, so the total buoyant force is
\[\vec F := -\int_{\partial\Omega}P\,\nu\,dS.\]
Substituting the pressure and applying the divergence theorem to each component,
\[F_i = \int_{\partial\Omega}\rho gx_3\nu_i\,dS
  = \rho g\int_{\partial\Omega}\langle x_3e_i,\nu\rangle\,dS
  = \rho g\int_\Omega\divg(x_3e_i)\,dV.\]
Since $\divg(x_3e_i) = \partial_i x_3$ equals $1$ if $i=3$ and $0$ otherwise, this gives
$F_1=F_2=0$ and
\[F_3 = \rho g\int_\Omega 1\,dV = \rho g\,\vol_3(\Omega).\]
The buoyant force is vertical, with magnitude $\rho g\vol_3(\Omega)$, which is exactly the weight
of the water displaced by $\Omega$, confirming Archimedes' principle. The divergence theorem
converts a surface integral of an unknown, position-dependent pressure into a volume integral
that only sees the shape of $\Omega$ through its volume.
\end{example}

% Source: Corsin Nick/Class Notes/Week 10.pdf, p. 12
Gauss' theorem is also a powerful tool in physics. Consider a fluid of density $\varrho$ moving
with velocity $\vec v$, and a surface element $d\vec A$ whose normal makes an angle $\theta$
with $\vec v$. The fluid crossing that element in time $\Delta t$ is exactly the fluid filling
the column of slant length $L = \Delta t\,\lvert\vec v\rvert$ erected on it along $\vec v$, and
that column stands at perpendicular height $L\cos\theta$. Its volume is therefore
\[V = \lvert d\vec A\rvert\,L\cos\theta = \Delta t\,\vec v\cdot d\vec A,\]
and the mass of fluid carried through the element is $\varrho V$.

\begin{ainote}
Corsin writes this as $V = \Delta t\,\vec v\cdot d\vec A\cdot\varrho$ and calls the result a
volume. The right-hand side is a mass: $\Delta t\,\vec v\cdot d\vec A$ is already the volume
swept, and the density only converts it. The two have been separated above, since the
distinction is needed immediately. The integral below is of $\varrho\vec v$ and therefore
measures mass per unit time, which is what the continuity equation balances against
$\partial\varrho/\partial t$.
\end{ainote}

% Generator: Opus 5 (high) --- FIG-W10-03, a redraw. The Sonnet 5 original was wrong in four
% ways, all of them visible only once rendered: it drew $\vec v$ twice in two different
% directions (a horizontal stream and a tilted arrow); the tilted arrow's label collided with
% the red one; it never drew the normal $d\vec A$ at all, so the $\theta$ arc sat between the
% velocity and nothing; and it labelled the height $h = \lvert d\vec A\rvert\cos\theta$, which
% is an area times a cosine.
%
% Geometry, all of it checked. The surface element is the segment (0,0)--(2,0); its normal
% dA runs along +y; the velocity makes the angle theta = 40 degrees with that normal, i.e.
% it runs at 50 degrees to the horizontal. The column swept in time Delta t is the element
% translated by L*(cos 50, sin 50) with L = 1.9, that is by (1.221, 1.455). Its perpendicular
% height is L*cos 40 = 1.455, which is exactly the y-extent of the parallelogram, so the h
% dimension line on the left measures the true height rather than an eyeballed one.
\begin{center}
\begin{tikzpicture}[>=Latex, scale=1.15]
  % The column swept through the element in time Delta t.
  \fill[purple!7] (0,0) -- (2,0) -- (3.221,1.455) -- (1.221,1.455) -- cycle;
  \draw[purple] (0,0) -- (1.221,1.455);
  \draw[purple, dashed] (1.221,1.455) -- (3.221,1.455);

  % The surface element itself, and the slant length of the column standing on it.
  \draw[purple, ultra thick] (0,0) -- (2,0);
  \node[purple, below, font=\scriptsize] at (1,-0.05) {$\lvert d\vec A\rvert$};
  \path (0,0) -- (1.221,1.455)
    node[midway, sloped, below, purple, font=\scriptsize] {$L=\Delta t\lvert\vec v\rvert$};

  % The normal, and the velocity at angle theta to it.
  \draw[blue, thick, ->] (2,0) -- (2,1.05) node[above, font=\scriptsize] {$d\vec A$};
  \draw[orange!90!black, thick, ->] (2,0) -- (3.221,1.455)
    node[right, font=\scriptsize] {$\Delta t\,\vec v$};
  \draw[green!50!black] (2,0.66) arc (90:50:0.66);
  \node[green!50!black, font=\scriptsize] at (2.17,0.42) {$\theta$};

  % The perpendicular height, drawn at the parallelogram's true y-extent.
  \draw[red!70!black, dashed] (-0.3,1.455) -- (1.221,1.455);
  \draw[red!70!black, dashed] (-0.3,0) -- (0,0);
  \draw[red!70!black, <->] (-0.3,0) -- (-0.3,1.455);
  \node[red!70!black, left, font=\scriptsize] at (-0.38,0.73) {$h=L\cos\theta$};
\end{tikzpicture}
\end{center}

% Source: Corsin Nick/Class Notes/Week 10.pdf, p. 12
Therefore, the fluid flowing out of any bounded $C^1$ domain $\Omega$ per unit time is
\[\int_{\partial\Omega} \varrho\vec v\cdot d\vec A = \int_\Omega \vec\nabla\cdot(\varrho
  \vec v)\,dx.\]
Of course, since the total amount of fluid is conserved, this is also the change in the
amount of fluid in $\Omega$ over time:
\[\int_\Omega \vec\nabla\cdot(\varrho\vec v)\,dx = -\frac{d}{dt}\int_\Omega \varrho\,dx
  = -\int_\Omega \frac{\partial\varrho}{\partial t}\,dx.\]
Since this holds for any $\Omega$, the \newterm{continuity equation}
\germanfor{Kontinuit\"atsgleichung} follows:
\[\vec\nabla\cdot(\varrho\vec v) + \frac{\partial\varrho}{\partial t} = 0.\]


% --- exercises relocated here from the dissolved sheets ---

% Originally: content/week-11.tex, \exercisesheet{11}, problem 11.5
% Source: exercises/Ex11_Analysis2_eng.pdf

\begin{exercise}[11.5 --- Radial vector fields \textnormal{(important)}]
\label{ex:11.5}
For $\alpha\in\mathbb{R}$ consider the vector field
\[V_\alpha(x) = \lvert x\rvert^\alpha x, \qquad x\in\mathbb{R}^n\setminus\{0\},\]
where $\lvert x\rvert$ is the usual Euclidean norm.
\begin{enumerate}[label=\textbf{\arabic*.}]
  \item Compute $\divg V_\alpha(x)$ for all $x\neq0$, and show that it vanishes identically
  for $\alpha=-n$.
  \item Using the Divergence theorem on $V_0$, show the following relation between the
  $n$-volume of a sphere and the $n-1$ volume of its boundary:
  \[r\,\vol_{n-1}(\partial B_r) = n\,\vol_n(B_r), \qquad \forall r>0.\]
  \item Show that
  \[\int_{\partial B_1} V_{-n}\cdot\nu\,d\vol_{n-1} > 0, \qquad \int_{B_1} \divg V_{-n}(x)\,dx
    = 0.\]
  Why doesn't this contradict the Divergence Theorem?
\end{enumerate}
\end{exercise}

\begin{aiexercise}[Flux of a cubic field through the sphere]
% Generator: Claude Opus 5 (High)
\label{ex:ai_flux_cubic_sphere}
Let $F(x,y,z) := (x^3,\,y^3,\,z^3)$ and let $\Omega := B_1(0) \subseteq \mathbb{R}^3$. Compute
the flux of $F$ out of $\partial\Omega$.
\begin{enumerate}[label=\textbf{\arabic*.}]
  \item Say in one sentence what the integrand of \cref{def:flux} would look like here if the
    sphere were parametrized directly, and why that route is unattractive.
  \item Compute $\divg F$, and observe that it depends on the point only through its distance
    $r$ from the origin.
  \item Evaluate the volume integral. The shell reading of polar coordinates in
    \cref{rem:polar_as_shells} does it without a Jacobian determinant.
\end{enumerate}
\end{aiexercise}

\begin{aiexercise}[An open surface has to be closed first]
% Generator: Claude Opus 5 (High)
\label{ex:ai_capping_a_paraboloid}
Let $S := \{(x,y,z) \in \mathbb{R}^3 : z = 1-x^2-y^2,\ z \geq 0\}$ be the piece of the
paraboloid lying above the $xy$-plane, oriented by the upward-pointing normal, and let
$F(x,y,z) := (x,\,y,\,z)$.
\begin{enumerate}[label=\textbf{\arabic*.}]
  \item \Cref{thm:gauss_divergence} applies to the boundary of a domain, and $S$ is not one.
    Which set $D$ must be adjoined to $S$ so that $S \cup D = \partial\Omega$ for a bounded
    domain $\Omega$, and what is the exterior unit normal along $D$?
  \item Compute $\int_{\partial\Omega}\langle F,\nu\rangle\,d\vol_2$ using the divergence
    theorem. You will need $\vol_3(\Omega)$, which is a one-line computation in cylindrical
    coordinates.
  \item Compute the flux through $D$ alone, directly from \cref{def:flux}, and deduce the flux
    of $F$ through $S$.
\end{enumerate}
\end{aiexercise}
